\lecture{12}
\title{Neural Networks}

\begin{document}

\subsection{Neurons}

Recall from \href{/6.790/lectures/11/#growth-functions-definition}{Lecture 11}
that the hypothesis class specified by regression models
does not contain a hypothesis that can correctly model
two-variable XOR;
the data pictured below is not linearly separable.
\begin{center}
    \frame{\includegraphics[width=0.1\textwidth]{lecture-11/linear-4}}
\end{center}
We need more general models than regression models.
In come \textbf{neural networks}.

\textbf{Definition.}
A single \textbf{neuron} is a function
$h \colon \mathbb{R}^d \to \mathbb{R}$,
made up of two parts:
\begin{itemize}
    \item The \textbf{pre-activation} of a neuron
    is $a(x) := w_0 + w^{\top} x$; that is,
    a weighted sum of $d$ inputs plus a bias $w_0$.
    \item The \textbf{activation / output} of a neuron
    is $h(x) := g(a(x))$, where $g \colon \mathbb{R} \to \mathbb{R}$
    is an \textbf{activation function}.
\end{itemize}
Some common choices of activation functions include:
\[ \text{(sigmoid $\sigma$)} ~~~ g(t) := \frac{1}{1 + e^{-t}}
~~~ \parallel ~~~
g(t) := \tanh(t) = \frac{e^t - e^{-t}}{e^t + e^{-t}}
~~~ \parallel ~~~
g(t) := \mathrm{sgn}(t)
~~~ \parallel ~~~
\text{(ReLU)} ~~~ g(t) := \max(0, t)  \]
\begin{remark}
    Note that $\tanh(x) = 2 \sigma(2x) - 1$,
    so the sigmoid and tanh activations
    are roughly the same.
\end{remark}

\subsection{Neural Networks: XOR Example}

Note that a single neuron with sigmoid activation
is identical to logistic regression.
This means that a single neuron
isn't enough to model two-variable XOR.
One fix, however, is to use a linear classifier
over a hardcoded \emph{feature} space.

\begin{center}
    \frame{\includegraphics[width=0.4\textwidth]{lecture-12/xor-diagram}}
\end{center}

As shown above, the following
$\mathbb{R}^2 \to \mathbb{R}^2$ feature map
allows us to compute $\mathrm{XOR}(x_1, x_2)$
via a linear classifier.
\[ \Phi \colon (x_1, x_2) \mapsto
\left(\mathrm{AND}(x_1, \overline{x}_2),
\mathrm{AND}(\overline{x}_1, x_2)\right). \]
In the diagram below, we represent
the process of first converting our inputs $(x_1, x_2)$
into features $\left(\mathrm{AND}(x_1, \overline{x}_2),
\mathrm{AND}(\overline{x}_1, x_2)\right)$,
then using the features as inputs to the output neuron.

\begin{center}
    \frame{\includegraphics[width=0.6\textwidth]{lecture-12/xor}}
\end{center}

Each input $(x_1, x_2)$ is a node,
each feature $\left(\mathrm{AND}(x_1, \overline{x}_2),
\mathrm{AND}(\overline{x}_1, x_2)\right)$ is a node,
and the output $\mathrm{XOR}(x_1, x_2)$ is a node.
Furthermore, an edge $u \to v$ is drawn if
$u$ is an input to the neuron whose output is $v$.

Here's the important point:
not only is the final output $\mathrm{XOR}(x_1, x_2)$ a neuron,
but the intermediate feature nodes
$\mathrm{AND}(x_1, \overline{x}_2)$ and
$\mathrm{AND}(\overline{x}_1, x_2)$
are neurons too!
Note that their activation function is $g(z) := \mathds{1}[z \geq 0]$.

The diagram above is an example of a \textbf{neural network}.
Let's generalize this to a definition.

\subsection{Neural Networks: Definition}

\textbf{Definition.}
A \textbf{neural network} consists of
one \textbf{input layer},
some number $L$ of \textbf{hidden layers},
and one \textbf{output layer}.
\begin{itemize}
    \item The \textbf{input layer}
    just reads in the input vector $x \in \mathbb{R}^d$.
    It's the identity map $h^{(0)} \colon \mathcal{X} \to \mathbb{R}^d$
    given by $h^{(0)}(x) := x$.
    \item The $k^{\text{th}}$ \textbf{hidden layer} is the map
    $h^{(k)} \colon \mathcal{X} \to \mathbb{R}^{N_k}$,
    whose output is a neuron activation
    of the output of the $(k - 1)^{\text{th}}$ layer.
    \[ h^{(k)}(x) = g\left(b^{(k)} + W^{(k)} h^{(k - 1)}(x)\right). \]
    Note that $N_k$ denotes
    the dimension of the output of $h^{(k)}$.
    Equivalently, $N_k$ is the number
    of neurons in the $k^{\text{th}}$ layer.

    By consequence, note that we must have
    $W^{(k)} \in \mathbb{R}^{N_k \times N_{k - 1}}$.
    \item The \textbf{output layer} is the map
    $h^{(L + 1)} \colon \mathcal{X} \to \mathbb{R}^{N_{L + 1}}$.
    It operates just like every hidden layer;
    its output is the final answer.
\end{itemize}
That's a whole bunch of jargon
to explain the relatively simple idea
of chaining $L$ layers of neurons together.

\begin{center}
    \frame{\includegraphics[width=0.3\textwidth]{lecture-12/neural-DAG}}
\end{center}

Here's how this looks in
the example of the XOR neural network.
Every single parameter in the neural network---that is,
every value of $b^{(k)}$ and
every entry of a matrix $W^{(k)}$---is
written in purple text.

\begin{center}
    \frame{\includegraphics[width=0.6\textwidth]{lecture-12/neural-xor}}
\end{center}

The upshot of neural networks
comes from their hidden layers.
Previously, we would just have to have \emph{known}, somehow, that
$\Phi(x) := \left(\mathrm{AND}(x_1, \overline{x}_2),
\mathrm{AND}(\overline{x}_1, x_2)\right)$ would be
a good set of features for encoding $\mathrm{XOR}(x_1, x_2)$.

With neural networks, however,
instead of \emph{hardcoding} the best features $\Phi(x)$
(or, in the neural network analogy, the best $h^{(L)}(x)$)
to use, we can instead \emph{learn} the best features
by optimizing the parameters $\{b^{(k)}, W^{(k)}\}_{k = 1}^{L + 1}$.

Hidden layers encode features.
That raises the question:
are hidden layers capable of encoding \emph{any} feature?

\subsection{Neural Network Universal Approximation}

The answer is yes, mostly.

\begin{theorem}{Universal Approximation}
    For every continuous $f \colon [0, 1]^d \to \mathbb{R}$
    and $\epsilon > 0$, there is a function $\hat{f}$
    that is a linear combination of ReLUs such that
    $\int_{[0, 1]^d} |f(x) - \hat{f}(x)| \, \mathrm{d}x < \epsilon$.
\end{theorem}

Furthermore, we can upper-bound
the complexity of such a function $\hat{f}$, somewhat.

\begin{theorem}{Universal Approximation Bound}
    Suppose $f \colon [0, 1]^d \to \mathbb{R}$
    is 1-Lipschitz continuous---that is,
    $|f(x) - f(y)| \leq \|x - y\|$
    for all $x, y \in [0, 1]^d$.
    Then for all $\epsilon > 0$, there is a function $\hat{f}$
    that is a linear combination of order $(1/\epsilon)^d$ ReLUs such that
    $\int_{[0, 1]^d} |f(x) - \hat{f}(x)| \, \mathrm{d}x < \epsilon$.
\end{theorem}

Interestingly, both of these theorems state
that only \emph{one} hidden layer is technically necessary---each
ReLU is a single node in that layer,
and the output layer is the linear combination of those nodes.

Still, though, as we'll see in future lectures,
it is often useful to use more than one hidden layer,
both from a computation perspective and a representation perspective.


\end{document}